\documentclass[a4paper,11pt]{article}

% Packages
\usepackage[utf8]{inputenc}
\usepackage{geometry}
\usepackage{amsmath, amssymb}
\usepackage{graphicx}
\usepackage{booktabs}
\usepackage{longtable}
\usepackage{siunitx}
\usepackage{enumitem}
\setlength{\marginparwidth}{2.5cm}
\usepackage[colorinlistoftodos]{todonotes}
\usepackage[hidelinks]{hyperref}

% Page setup
\geometry{margin=2.5cm}

% Convenience macros
\newcommand{\code}[1]{\texttt{#1}}
\newcommand{\TODOfill}{\todo[inline]{Fill in.}}

% Title info
\title{Assignment 4: Self-Supervised Representation Learning\\ for Time Series Classification}
\author{Frederik Appel Vardinghus-Nielsen}
\date{\today}

\begin{document}

\maketitle
\listoftodos
\tableofcontents
\newpage

%======================================================================
\section{Part A: Understand the methods and connect paper to code}
%======================================================================

\subsection{The three learning ideas}

\subsubsection{T-Loss (Franceschi et al., NeurIPS 2019)}
The general idea is that sub-series of a reference series (which may be a sub-series of a larger time series) is a positiv sample, which sub-series of another series are negative samples. They use a CNN with dilation for increased receptive field. The series are left padded with zero if they are too short for the dilation of the network. If the series are long, more time stamps have representations in the end which is fixed with a max pool.

\begin{itemize}
\item Positive/negative samples: a positive example $x^{pos}$ is a sub-series of the reference series $x^{ref}$, which itself is randomly sampled in length and location. A sub-series of another time series is negative.
\item Objective: triple loss with one anchor, one positive, and $K\geq1$ negative samples $x_{k}^{neg}$. The objective is to have the distance between $x^{pos}$ and $x^{ref}$ be smaller than the sum of the distances between $x^{ref}$ and $x_k^{neg}$.
\item Fixed length problem: the convolutional network ends with a representation $Y\in\mathbb{R}^{160 \times L}$ which is max pooled into $z\in\mathbb{R}^{160}$ and thereafter transformed through $f(x)=Wz+b$ where $W\in\mathbb{R}^{320\times160}$. The final dot product are therefore always between vectors of length 320. If the time series chosen in the beginning are too short for the full CNN they are left zero padded.
\item Classification: an SVM with a radial basis function is trained on the features with their labels now avalable. It is evaluated in a train/test split.
\end{itemize}

\subsubsection{TS2Vec (Yue et al., AAAI 2022)}
The general idea is to generate related views with masking and random overlapping segments of the representations. The hierarchical loss furthermore lets the encoder learn representations which take into account a spectrum of resolutions (granularity of the series).

\begin{itemize}
\item Positive/negative samples: two different views of the same time series instance (segment) are achieved through random cropping and timestamp masking. The original time series is cropped to produced two distinct, overlapping time series. After and input projection layer, these are both randomly masked at specific time stamps (single timestamp, all dimensions for multivariate TS) by setting values to zero.
\item Objective: Dual loss, the sum of two distinct expressions:
\begin{itemize}
\item Temporal contrastive loss which compares views of the same overlapping series at the same and different timestamps. Similarity at same timestamp between two different views is positive, self-similarity (same view, no masking) at different timestamps is negative.
\item Instance-wise contrastive loss compares to other time series. Positives are two views of the same time series at the same timestamp. Negatives are both view from another time series at the same time stamp.
\end{itemize}
The loss is calculated recursively at larger and larger granularity with a dilation of 2. The first level compares every single timestamp with all others. The second max pools with kernel size 2 and this continues until the full series are used.
\item Fixed length problem: only the overlapping timestamps are used after passing the views through the encoder. The loss loops over time stamps and gives a single loss for them all.
\item Classification: they use instance-level representations (the whole segment of time series, not just the overlap). Instances are passed through the encoder and a max pooling is performed over all timestamps, giving a feature vector of dimension $K$, which is the number of representation dimensions in the last layer. SVM is trained on this $K$-dimensional vector for classification.
\end{itemize}

\subsubsection{TF-C (Zhang et al., NeurIPS 2022)}
They seek to enforce discrimination in different domains
This encourages that the time domain is robust to augmentation, the frequency domain is robust to augmentation, and the time-frequency domain enforces temporal and spectral consistency.

\begin{itemize}
\item Positive/negative samples: 
\begin{itemize}
\item Time: (1) Encodings of a time series and its augmented (jittering, scaling, time-shifts, neighborhood segments) view are similar and (2) encodings of two different time series (no augmentations) are dissimilar.
\item Frequency: (1) Encodings of the Fourier transform of a time series and its agumented view are similar and (2) encodings of two different series (no augmentations) are dissimilar.
\item Time-frequency: (1) Time series and their Fourier transform must be similar after first encoding into respective domains and thereafter encoding into common time-frequency domain, (2), augmented views are dissimilar.
\end{itemize}
\item Objective: Combined loss in time-, frequency-, and time-frequency domain.
\item Fixed length problem: The training datasets all contain series of same length. When using the model on a tuning dataset and testing, the series are left zero padded. 
\item Classification: the full architecture of $z_i^T=R_T(G_T(x_i^T))$ and $z_i^F=R_F(G_F(_i^F))$ is used and concatenated into $[z_i^T, z_i^F]$ and this is passed through a 2-layer fully ocnnected network which projects into the number of classes needed. This classifier is jointly trained with the full pipeline on the target dataset as fine tuning.
\end{itemize}


\subsection{Paper-to-code map}

For each element: (i) closest paper location, (ii) inputs and outputs, (iii) immediate
caller or call site, and (iv) role in training or representation extraction.
Verified checkout commits: T-Loss \code{4aff592}, TS2Vec \code{b0088e1}, TF-C \code{9667582}.

\subsubsection{T-Loss}

\paragraph{\code{losses/triplet\_loss.py: TripletLoss.forward}}
\begin{enumerate}[label=(\roman*)]
    \item Eq. (1)
    \item
    \textbf{Input:} tensor \texttt{(B, C, L)}, where \texttt{B} is the batch size, \texttt{C} is the number of channels, and \texttt{L} is the length of the time series (equal length case).\\
    \textbf{Output:} Loss as a scalar.
    \item \texttt{TimeSeriesEncoderClassifier.fit\_encoder}
    \item It is the loss which dictates which time series are supposed to have similar and dissimilar representations.
\end{enumerate}

\paragraph{\code{networks/causal\_cnn.py: CausalConvolutionBlock.forward}}
\begin{enumerate}[label=(\roman*)]
\item Figure 2(b) in the paper.
\item 
    \textbf{Input:} tensor \texttt{(B, C, L)}, where \texttt{B} is the batch size, \texttt{C} is the number of channels, and \texttt{L} is the length of the time series (equal length case).\\
    \textbf{Output:} tensor \texttt{(B, C, L)}, where \texttt{B} is the batch size, \texttt{C} is the number of channels, and \texttt{L} is the length of the time series (equal length case). This is the results after passing through the convolutional block.
\item \texttt{CausalCNN} where is is called for each layer in the full network.
\item It's the basic block which dictates the neural network achitecture and therefore the inductive bias of the representation learning.
\end{enumerate}

\paragraph{\code{networks/causal\_cnn.py: CausalCNNEncoder.forward}}

\begin{enumerate}[label=(\roman*)]
\item Section 4 describes the full architecture. Figure 2 illustrates the convolutional blocks and dilation which are repeated in the first part of the encoder.
\item 
    \textbf{Input:} tensor \texttt{(B, C, L)}, where \texttt{B} is the batch size, \texttt{C} is the number of channels, and \texttt{L} is the length of the time series (equal length case).\\
    \textbf{Output:} tensor \texttt{(B, C)}, where \texttt{B} is the batch size and \texttt{C} is the number of channels. The time dimension $L$ is removed with max pooling squeezing the time dimension.
\item \texttt{TimeSeriesEncoderClassifier.encode}
\item It wraps the whole classification pipeline of the encoder and calls it on the passed batch.
\end{enumerate}

\paragraph{\code{scikit\_wrappers.py: TimeSeriesEncoderClassifier.fit\_encoder}}
\begin{enumerate}[label=(\roman*)]
\item Section 4 on architecture and Section S2 on hyperparameters.
\item 
    \textbf{Input:} Training set $X$.\\
    \textbf{Output:} The fitted encoder.
\item \texttt{TimeSeriesEncoderClassifier.fit}
\item It performs unsupervised/self-supervised fitting of the encoder.
\end{enumerate}

\paragraph{\code{scikit\_wrappers.py: TimeSeriesEncoderClassifier.encode}}
\begin{enumerate}[label=(\roman*)]
\item Paragraph 4 and 5 in Section 4 on encoder architecture.
\item 
    \textbf{Input:} The test set $X$. \\
    \textbf{Output:} Features of $X$.
\item Depends on the context. In general \texttt{fit()} calls it during training to sue the final encoder to produce features before training SVM, while \texttt{score()} calls is during testing to produce features for SVM classification.
\begin{itemize}
\item \texttt{fit\_encoder()} calls it after each training epoch if certain conditions for early stopping are present. This doesn't make sense for fully self-supervised learning, so labels must be 
provided.
\item \texttt{fit()} calls it to encode the training set once training is done such that SVM can be training on the training features.
\item \texttt{encode\_window()} calls it to encode a test set on a sliding window. Useful if temporal features are needed for regression/forecasting.
\item \texttt{predict()} and \texttt{score()} call it to encode a set and then perform prediction or calculate score of classification. Mostly used for testing.
\end{itemize}
\item It uses the final encoder (unless used during early stopping) to encode features from either train or test sets before classification or prediction.
\end{enumerate}

\paragraph{Required specifics}
\begin{itemize}
    \item How \code{fit\_encoder} selects the fixed-length or varying-length loss: It checks wether a \texttt{numpy} sum over the training set yields \texttt{nan}s in some entries, which then indicates unequal lengths.
    \item How \code{TripletLoss.forward} samples reference, positive and negative intervals: 
    \begin{itemize}
    \item Reference/anchor: each time series is used as reference once per epoch. The length is chosen at random between the length of the positive example and the length of the full series. The position is chosen to overlap the positive example.
    \item Negative: a specified number of random time series are chosen as negatives for each anchor. The anchor itself can contain a negative sample.
    \item Positive: chosen as a random sub-series of each anchor. In practice the length and position of the positive is chosen first.
    \end{itemize}
    \item How \code{CausalCNNEncoder} produces one fixed-length vector: it uses AdaptiveMaxPool1d to max pool over the time dimension. This produces a feature vector which only has length one in the remaining "time" dimension.
    \item How to call \code{fit\_encoder} without labels: call it without passing training labels \texttt{y}. It defaults to none and performs unsupervised learning without early stopping.
\end{itemize}

\subsubsection{TS2Vec}

\paragraph{\code{ts2vec.py: TS2Vec.fit}}
\begin{enumerate}[label=(\roman*)]
\item Section 2.3 on Contextual Consistency which describes the timestamp masking and random cropping which serve as the augmentations for robustness, and Section 2.4 on Hierarchical Contrasting which is described in Algorithm 1. The whole loss is shown in Equations (1), (2), and (3). 
\item 
    \textbf{Input:} Train data, number of epochs to train, maximum number of iterations (batches) to train, and whether to print training loss after each epoch. \\
    \textbf{Output:} Log of losses during training.
\item It's called in \texttt{train.py}.
\item It contains the central training loop which includes cropping, augmentation, and loss.
\end{enumerate}

\paragraph{\code{ts2vec.py: TS2Vec.\_eval\_with\_pooling}}
\begin{enumerate}[label=(\roman*)]
\item Sections 3.1, 3.2, and 3.3 on applications. The evaluation/representation of the time series varies depending on the downstream task.
\item 
    \textbf{Input:} Test data \texttt{x} on which to evaluate the averaged model, whether to use a mask, whether to slice.\\
    \textbf{Output:} Representations from running the encoder and the chosen pooling strategy.
\item \texttt{TS2Vec.fit()}
\item It applies the trained model to a test set and the pooling strategy which is chosen based upon the downstream task.
\end{enumerate}

\paragraph{\code{ts2vec.py: TS2Vec.encode}}
\begin{enumerate}[label=(\roman*)]
\item Section 2.2 on Model Architecture and Sections 3.1, 3.2, and 3.3 on downstream tasks.
\item 
    \textbf{Input:} Test data to calculate representations from. Masking, sliding length and padding for representation altering, encoding window for restriction before encoding, causal, and batch sizes are additional inputs. \\
    \textbf{Output:} The representations from encoding the test data.
\item \texttt{anomaly\_detection.py}, \texttt{classification.py}, and \texttt{forecasting.py} for the three tasks, respectively.
\item Uses the fitted encoder to encode test data for use in downstream task.
\end{enumerate}

\paragraph{\code{models/losses.py: hierarchical\_contrastive\_loss}}
\begin{enumerate}[label=(\roman*)]
\item Equations (1), (2), and (3).
\item 
    \textbf{Input:} Two different reprensentation of the same batch to calculate loss between. \\
    \textbf{Output:} Loss between the representations as defined by Equation (3).
\item \texttt{TS2Vec.fit()}
\item Calculates loss which dictates training.
\end{enumerate}

\paragraph{\code{models/losses.py: instance\_contrastive\_loss}}
\begin{enumerate}[label=(\roman*)]
\item Equation (2).
\item
    \textbf{Input:} Two representations of the same batch after encoding during training. \\
    \textbf{Output:} The instance contrastive loss between the two representations. Two representations at the same timestamp are positive, while same representation of two different time series are negative.
\item \texttt{hierarchical\_contrastive\_loss}.
\item Forces the model to have representations of the same series in different contexts and with different masks to be similar, while keeping reprensentations of other series dissimilar.
\end{enumerate}

\paragraph{\code{models/losses.py: temporal\_contrastive\_loss}}
\begin{enumerate}[label=(\roman*)]
\item Equation (2).
\item 
    \textbf{Input:} Two representations of the same batch after encoding during training.\\
    \textbf{Output:} The temporal contrastive loss between the two reprensentations. The same timestamp for the two representations is a positive pair, while different timestamps in both the same and the other representation are negative pairs.
\item \texttt{hierarchical\_contrastive\_loss}.
\item Encourages the model to discriminate representations over time. Representations at the same timestamps should have similar representations while temporal differences should be reflected in different representations.
\end{enumerate}

\paragraph{\code{models/encoder.py: TSEncoder.forward}}
\begin{enumerate}[label=(\roman*)]
\item Section 2.2, paragraphs 2 and 3.
\item 
    \textbf{Input:} Data to be encoded. It has been cropped and augmented before. \\
    \textbf{Output:} Encoded batch.
\item \texttt{TS2Vec.fit()} and \texttt{\_eval\_with\_pooling}.
\item Contains the convolutional part of the model architecture.
\end{enumerate}

\paragraph{\code{models/encoder.py: generate\_binomial\_mask}}
\begin{enumerate}[label=(\roman*)]
\item The paragraph \textbf{Timestamp Masking} in Section 2.3.
\item 
    \textbf{Input:} \texttt{B}, \texttt{T}, and \texttt{p}; batch size, time length, and masking probability, respectively.\\
    \textbf{Output:} A boolean tensor with shape \texttt{(B, T)} with values \texttt{False} and \texttt{True} indicating masking or not, respectively.
\item \texttt{TSEncoder.forward()}
\item Performs the timestamp masking which forces robustness and context understanding of the model.
\end{enumerate}

\paragraph{\code{models/encoder.py: generate\_continuous\_mask}}
\begin{enumerate}[label=(\roman*)]
\item Not mentioned in article.
\item 
    \textbf{Input:} \texttt{B}, \texttt{T}, \texttt{n}, and \texttt{l}; batch size, time length, number of masked sequences, and length of masked sequences as a fraction of full length, respectively.\\
    \textbf{Output:} A boolean tensor with shape \texttt{(B, T)} with values \texttt{False} and \texttt{True} indicating masking or not, respectively. The \texttt{while} loop in \texttt{hierarchical\_contrastive\_loss} matches the loop in Algorithm 1.
\item \texttt{TSEncoder.forward()}
\item Creates continuous runs of masking. They may overlap.
\end{enumerate}

\paragraph{Required specifics}
\begin{itemize}
    \item Overlapping crops in \code{TS2Vec.fit}: They are created by in sequence determining the length of the overlap of the crops, the start and end of the overlap, the extension of crops 1 to the left, and finally the extension of crop 2 to the right.
    \item Where timestamp masking is applied: In TSEncoder.forward(). This is only applied during training.
    \item Mapping of the two elementary losses and hierarchical pooling to Equations 1--3 and Algorithm 1: \texttt{temporal\_contrastive\_loss}, \texttt{instance\_contrastive\_loss}, and \texttt{hierarchical\_contastive\_loss} match Equations (1)--(3), respectively.
    \item The \code{TS2Vec.encode} option returning one vector per complete series: Set the argument \texttt{encoding\_window="full\_series"}. This causes a max pool of the features over all time stamps after encoding.
\end{itemize}

\subsubsection{TF-C}

\paragraph{\code{code/TFC/dataloader.py: Load\_Dataset.\_\_init\_\_}}
\begin{enumerate}[label=(\roman*)]
\item Sections 4.1 and 4.2 on data augmentation in time and frequency domains.
\item 
    \textbf{Input:} A dataset with time series and labels, a config, training mode, max size of a debugging dataset, and whether to debug (use a smaller dataset). \\
    \textbf{Output:} No output. Saves two version of the time series: the augmented version and the Fourier transformed and then agumented.
\item Three times in \texttt{data\_generator}; once for each of training, fine-tuning, and testing.
\item Performs the steps of loading and augmenting data before it is passed through the first set of encoders.
\end{enumerate}

\paragraph{\code{code/TFC/dataloader.py: Load\_Dataset.\_\_getitem\_\_}}
\begin{enumerate}[label=(\roman*)]
\item Has no direct counterpart in the article. It just serves to pop out a single time series, its augmentations, and transforms.
\item 
    \textbf{Input:} Index of the requested time series.\\
    \textbf{Output:} The requested time series, its augmentation in time, the Fourier transform, its augmentation in frequency, and the labels. if run at fine tuning or test time, the augmentations are the same as the original representations.
\item \texttt{model\_pretrain}, \texttt{mode\_finetune}, and \texttt{mode\_test}.
\item Is a wrapper for providing the augmented and transformed views and labels during training, fine tuning, and test.
\end{enumerate}

\paragraph{\code{code/TFC/dataloader.py: data\_generator}}
\begin{enumerate}[label=(\roman*)]
\item No direct counterpart. Paragraph 4 of Section 1 explains the concept of pre-training, fine-tuning, and then testing. This split is done by this generator. 
\item 
    \textbf{Input:} Path to the source data, path to the target (fine tuning and test dataset), configs, training mode (pre-train, fine tuning, or test), and whether to run in debugging mode.\\
    \textbf{Output:} Loaders for the three different cases and datasets: training, fine tuning, and testing.
\item \texttt{main.py}
\item Helps loading the correct datasets with the corresponding augmentations for pre-training, fine tuning, and testing.
\end{enumerate}

\paragraph{\code{code/TFC/dataloader.py: the \code{fft.fft(...).abs()} operation}}
\begin{enumerate}[label=(\roman*)]
\item Section 4.2, paragraph 1, where the frequency representation is described as the Fourier transformation of the time series.
\item 
    \textbf{Input:} The whole dataset of time series. \\
    \textbf{Output:} The absolute values of the Fourier transform of the time series.
\item It is saved in the class instance as the spectral representation of the time series. During training, it is called in \texttt{TFC.forward()}.
\item Is responsible for the spectral representation which the authors deem valuable for time series representation.
\end{enumerate}

\paragraph{\code{code/TFC/augmentations.py: DataTransform\_TD}, \code{DataTransform\_FD}, \code{remove\_frequency}, \code{add\_frequency}}
\begin{enumerate}[label=(\roman*)]
\item Sectio 4.2, paragraph 2 on augmentations of the Fourier transform of the time series.
\item 
    \textbf{Input:} Fourier transform of a single time series. \\
    \textbf{Output:} An augmented Fourier transform where frequencies are added and removed at random.
\item \texttt{Load\_Dataset.\_\_init\_\_()} when augmenting the Fourier transformed time series.
\item Performs augmentations of the absolute Fourier transform of the raw time series in order to provide positive pairs in the spectral representations. The code is wrongly implemented and does not actually remove frequencies but rather in practice amplifies others.
\end{enumerate}

\paragraph{\code{code/TFC/model.py: TFC.\_\_init\_\_}, \code{TFC.forward}, \code{target\_classifier.forward}}
\begin{enumerate}[label=(\roman*)]
\item Calls and structures the overall architecture and data flow in the model. \texttt{init} defines the four encoders, \texttt{forward} defines the four representations when the time series, the Fourier transforms, and their augmentments are passed through the encoders, and \texttt{target\_classifier.forward} is simply a small, fully connected neural network which is training during fine tuning and used for classification during testing.
\item 
    \textbf{Input:} Training data and its Fourier transform.\\
    \textbf{Output:} The four representations encoded by the various combinations of the encoders.
\item \texttt{main.py}
\item This structures all of the defined architectures into the whole model and defines the flow of data through the various networks and classifiers.
\end{enumerate}

\paragraph{\code{code/TFC/loss.py: NTXentLoss\_poly.forward}}
\begin{enumerate}[label=(\roman*)]
\item Equations (1), (2), and (3).
\item 
    \textbf{Input:} Representations $Z$ and $\tilde{Z}$.\\
    \textbf{Output:} Total loss from postive and negative samples in the two representations.
\item It is called in \texttt{trainer.py: model\_pretrain} and \texttt{model\_finetune} for use during training and fine tuning.
\item Implements the well known NT-Xent loss between the various positive and negative samples in the model.
\end{enumerate}

\paragraph{\code{code/TFC/trainer.py: Trainer}, \code{model\_pretrain}, \code{model\_finetune}, \code{model\_test}}
\begin{enumerate}[label=(\roman*)]
\item Sections 5.1 and 5.2 on results from pre-training and fine-tuning.
\item 
    \textbf{Input:} Model, optimizer, classifier, classifier optimizer, dataloaders, and configuration files.\\
    \textbf{Output:} No returns. Prints logging during training or fine tuning and saves models and results.
\item \code{main.py}
\item \code{Trainer} splits the model in pre-training and fine-tuning and runs the overall logic for pre-training and fine-tuning and saves the models, \code{mode\_pretrain} runs the pre training pipeline and returns the loss, and \code{model\_finetune} does the same for the test datasets and reports statistics on the results.
\end{enumerate}

\paragraph{\code{code/TFC/main.py: the \code{subset} setting, call to \code{data\_generator}, pretrained-checkpoint loading}}
\begin{enumerate}[label=(\roman*)]
\item No directly corresponding parts of the article.
\item 
    \textbf{Input:} Not a function, so no inputs. However, needs configuration files and datasets.\\
    \textbf{Output:} No outputs but logs progress during pre-training of fine-tuning. The models which are trained are saved to the local machine during this function.
\item No immediate caller.
\item Top level file for entry for both pre-training and fine-tuning.
\end{enumerate}

\subsubsection{Following one TF-C batch through the pipeline}
\begin{enumerate}
\item Data is loaded by \code{Load\_Dataset} where the FFT and augmentations are created.
\item \code{data\_generator} calls the loader and creates training datasets.
\item The loop in \code{model\_pretrain} loads a batch.
\item The batch is passed through the encoders.
\item Loss is calculated from the representations of the batch.
\item The loss is used for backpropagation.
\end{enumerate}

\paragraph{FFT dimension, spectrum conversion, \code{fft} vs.\ \code{rfft}}
\begin{itemize}
\item \textbf{FFT dimension:} over the time dimension, two sided to keep matching length with the time domain part.
\item \textbf{Spectrum conversion:} absolute value/magnitude is used.
\item \code{fft} or \code{rfft}: \code{fft} is used. This keeps the same length as the original time series.
\end{itemize}

\paragraph{Active time-domain augmentation; the two frequency operations and how they are combined}
\begin{itemize}
\item \textbf{Time domain:} the paper mentions a bank of multiple different of which a radnom should be chosen. In practice only jitter is used.
\item \textbf{Frequency domain:} Frequencies (bins) are added and removed at random. The two augmentation are added which results in the removed frequencies not actually being removed, more halved.
\end{itemize}

\paragraph{Attributes defining the time encoder, frequency encoder and the two projectors; the four outputs of \code{TFC.forward} in terms of $h$ and $z$}
\begin{itemize}
\item \textbf{Time and frequency encoders (same parameters):} length of the time series, dimension of the intermediate layer in the transformer, and number of heads of the attention mechanism.
\item \textbf{Projectors (same parameters):} Dimension of the first linear layer, dimension of batch normalisation, and dimensions of last linear layer.
\item \textbf{Relation of outputs to paper notation:}
\begin{itemize}
\item \code{h\_time} corresponds to $\mathbf{h}^T$ which is the encoding of the time series by the time encoder, utilizing transformer architecture.
\item \code{z\_time} corresponds to $\mathbf{z}^T$ which is the projection of $\mathbf{h}^T$ by the time projector, utilizing linear layers and ReLU.
\item \code{h\_freq} corresponds to $\mathbf{h}^F$ which is the encoding of the Fourier transofmr of the time series by the frequency encoder, utilizing transformer architecture.
\item \code{z\_freq} corresponds to $\mathbf{z}^F$ which is the projection of $\mathbf{h}^F$ by the frequency projector, utilizing linear layers and ReLU.
\end{itemize}
\end{itemize}

\paragraph{Loss terms in \code{model\_pretrain} vs.\ Equations 1--4}
\begin{itemize}
\item \code{loss_t = nt_xent_criterion(h_t, h_t_aug)} is the within-time loss.
\item \code{loss_f = nt_xent_criterion(h_f, h_f_aug)} is the within-frequency loss.
\end{itemize}

\paragraph{Tensors concatenated for the downstream classifier; is the encoder frozen or fine-tuned?}
\TODOfill

\paragraph{\code{main.py} $\rightarrow$ \code{data\_generator} $\rightarrow$ \code{Trainer} for FD-A $\rightarrow$ FD-B}
Number of source and target-training samples actually used, whether \code{FD\_B/val.pt} is loaded,
how often the test loader is evaluated, and whether test results influence the reported best result.
Discrepancies from the paper's stated protocol are recorded, not repaired.
\TODOfill

\paragraph{Device-specific tensor construction in \code{remove\_frequency} and \code{add\_frequency}}
\TODOfill

%======================================================================
\section{Part B: Raw-signal baselines}
%======================================================================

\subsection{Protocol}
Datasets: FordA, ArticularyWordRecognition, Epilepsy. Classifiers: logistic regression and
RBF-kernel SVM on the vectorised raw signal. Label regimes: 100\% and a stratified 10\% subset.
Three fixed seeds; the 10\% subset of a given seed is reused by every method.
Hyperparameters selected on training data only.
\TODOfill

\subsection{Results}

\begin{table}[htbp]
    \centering
    \small
    \begin{tabular}{l l l c c}
        \toprule
        Dataset & Model & Labels & Accuracy & Macro-F1 \\
        \midrule
        FordA & Logistic regression & 100\% & & \\
        FordA & Logistic regression & 10\%  & & \\
        FordA & RBF-SVM             & 100\% & & \\
        FordA & RBF-SVM             & 10\%  & & \\
        \midrule
        AWR & Logistic regression & 100\% & & \\
        AWR & Logistic regression & 10\%  & & \\
        AWR & RBF-SVM             & 100\% & & \\
        AWR & RBF-SVM             & 10\%  & & \\
        \midrule
        Epilepsy & Logistic regression & 100\% & & \\
        Epilepsy & Logistic regression & 10\%  & & \\
        Epilepsy & RBF-SVM             & 100\% & & \\
        Epilepsy & RBF-SVM             & 10\%  & & \\
        \bottomrule
    \end{tabular}
    \caption{Raw-signal baselines. Test accuracy and macro-F1 as mean $\pm$ standard deviation over three paired seeds.}
    \label{tab:baselines}
\end{table}

\subsection{Comments}
\TODOfill

%======================================================================
\section{Part C: Pretrain and evaluate T-Loss and TS2Vec}
%======================================================================

\subsection{Protocol}
Per dataset and seed: pretrain on the unlabelled official training series, extract one
fixed-length representation per series, freeze the encoder, fit the probe on the 10\% and
100\% labelled training representations, and evaluate once on the official test set.
\TODOfill

\paragraph{Principal representation test}
The probe treated as the principal representation test is: \TODOfill

\subsection{Results}

\begin{table}[htbp]
    \centering
    \small
    \begin{tabular}{l l l l c c}
        \toprule
        Dataset & Method & Probe & Labels & Accuracy & Macro-F1 \\
        \midrule
        FordA & T-Loss & Logistic regression & 100\% & & \\
        FordA & T-Loss & Logistic regression & 10\%  & & \\
        FordA & T-Loss & RBF-SVM             & 100\% & & \\
        FordA & T-Loss & RBF-SVM             & 10\%  & & \\
        FordA & TS2Vec & Logistic regression & 100\% & & \\
        FordA & TS2Vec & Logistic regression & 10\%  & & \\
        FordA & TS2Vec & RBF-SVM             & 100\% & & \\
        FordA & TS2Vec & RBF-SVM             & 10\%  & & \\
        \midrule
        AWR & T-Loss & Logistic regression & 100\% & & \\
        AWR & T-Loss & Logistic regression & 10\%  & & \\
        AWR & T-Loss & RBF-SVM             & 100\% & & \\
        AWR & T-Loss & RBF-SVM             & 10\%  & & \\
        AWR & TS2Vec & Logistic regression & 100\% & & \\
        AWR & TS2Vec & Logistic regression & 10\%  & & \\
        AWR & TS2Vec & RBF-SVM             & 100\% & & \\
        AWR & TS2Vec & RBF-SVM             & 10\%  & & \\
        \midrule
        Epilepsy & T-Loss & Logistic regression & 100\% & & \\
        Epilepsy & T-Loss & Logistic regression & 10\%  & & \\
        Epilepsy & T-Loss & RBF-SVM             & 100\% & & \\
        Epilepsy & T-Loss & RBF-SVM             & 10\%  & & \\
        Epilepsy & TS2Vec & Logistic regression & 100\% & & \\
        Epilepsy & TS2Vec & Logistic regression & 10\%  & & \\
        Epilepsy & TS2Vec & RBF-SVM             & 100\% & & \\
        Epilepsy & TS2Vec & RBF-SVM             & 10\%  & & \\
        \bottomrule
    \end{tabular}
    \caption{Frozen-representation classification. Mean $\pm$ standard deviation over three paired seeds.}
    \label{tab:ssl-clean}
\end{table}

\subsection{Comparison with the raw-signal baselines}
\TODOfill

\subsection{Training time and representation dimensionality}

\begin{table}[htbp]
    \centering
    \small
    \begin{tabular}{l l c c}
        \toprule
        Dataset & Method & Pretraining time & Representation dimension \\
        \midrule
        FordA & T-Loss & & \\
        FordA & TS2Vec & & \\
        AWR & T-Loss & & \\
        AWR & TS2Vec & & \\
        Epilepsy & T-Loss & & \\
        Epilepsy & TS2Vec & & \\
        \bottomrule
    \end{tabular}
    \caption{Training time and representation dimensionality.}
    \label{tab:cost}
\end{table}

%======================================================================
\section{Part D: Robustness to contiguous sensor dropout}
%======================================================================

\subsection{Corruption design}
Dropout proportion, number of intervals and replacement convention, fixed before inspecting
any result, plus a brief justification. The same corrupted inputs are applied to all methods.
\TODOfill

\subsection{Results (100\%-label models)}

\begin{table}[htbp]
    \centering
    \small
    \begin{tabular}{l l c c c c}
        \toprule
        & & \multicolumn{2}{c}{Accuracy} & \multicolumn{2}{c}{Macro-F1} \\
        \cmidrule(lr){3-4} \cmidrule(lr){5-6}
        Dataset & Method & Clean & Corrupted ($\Delta$) & Clean & Corrupted ($\Delta$) \\
        \midrule
        FordA & Logistic regression (raw) & & & & \\
        FordA & RBF-SVM (raw)             & & & & \\
        FordA & T-Loss                    & & & & \\
        FordA & TS2Vec                    & & & & \\
        \midrule
        AWR & Logistic regression (raw) & & & & \\
        AWR & RBF-SVM (raw)             & & & & \\
        AWR & T-Loss                    & & & & \\
        AWR & TS2Vec                    & & & & \\
        \midrule
        Epilepsy & Logistic regression (raw) & & & & \\
        Epilepsy & RBF-SVM (raw)             & & & & \\
        Epilepsy & T-Loss                    & & & & \\
        Epilepsy & TS2Vec                    & & & & \\
        \bottomrule
    \end{tabular}
    \caption{Clean versus contiguous-dropout test performance for the 100\%-label models.}
    \label{tab:dropout}
\end{table}

\subsection{Discussion}
Whether any apparent robustness comes from the representation, the classifier or the
preprocessing convention.
\TODOfill

%======================================================================
\section{Part E: Focused TS2Vec continuous masking ablation}
%======================================================================

\subsection{Setup}
On FordA: standard (binomial) TS2Vec masking versus continuous masking, with encoder,
training budget, downstream probe, label regime and the three seeds matched.
\TODOfill

\subsection{Results}

\begin{table}[htbp]
    \centering
    \small
    \begin{tabular}{l c c c c}
        \toprule
        & \multicolumn{2}{c}{Accuracy} & \multicolumn{2}{c}{Macro-F1} \\
        \cmidrule(lr){2-3} \cmidrule(lr){4-5}
        Masking & Clean & Corrupted ($\Delta$) & Clean & Corrupted ($\Delta$) \\
        \midrule
        Binomial (default) & & & & \\
        Continuous         & & & & \\
        \bottomrule
    \end{tabular}
    \caption{FordA: standard versus continuous TS2Vec masking, clean and contiguous-dropout test sets.}
    \label{tab:masking-ablation}
\end{table}

\subsection{Interpretation}
Does continuous masking change (i) clean performance, (ii) corrupted performance, and
(iii) sensitivity to corruption relative to the model's own clean score?
\TODOfill

%======================================================================
\section{Part F: TF-C for industrial bearing fault classification (\texorpdfstring{FD-A $\rightarrow$ FD-B}{FD-A to FD-B})}
%======================================================================

\subsection{Execution record}

\begin{table}[htbp]
    \centering
    \small
    \begin{tabular}{l p{9.5cm}}
        \toprule
        Item & Value \\
        \midrule
        Commands (pretraining) & \\
        Commands (fine-tuning / testing) & \\
        Environment (Python, PyTorch, OS) & \\
        Hardware & \\
        Seed & \\
        Data actually consumed & \\
        Training time & \\
        Checkpoint used & \\
        \bottomrule
    \end{tabular}
    \caption{TF-C FD-A $\rightarrow$ FD-B execution record.}
    \label{tab:tfc-record}
\end{table}

\subsection{Classification result}

\begin{table}[htbp]
    \centering
    \small
    \begin{tabular}{l c c}
        \toprule
        Run & Accuracy & Macro-F1 \\
        \midrule
        This run (FD-B test) & & \\
        TF-C supplement, Table 5 & & \\
        \bottomrule
    \end{tabular}
    \caption{FD-B classification result, compared cautiously with Table 5 of the TF-C supplement. Exact numerical reproduction is not expected.}
    \label{tab:tfc-result}
\end{table}

\subsection{What was transferred, where labels entered, and encoder status}
What was transferred from FD-A, where FD-B labels entered the pipeline, and whether the final
encoder was frozen or updated.
\TODOfill

\subsection{Material differences between the paper's protocol and the supplied execution path}
Subset use, loss composition, validation use and test-set access, based on the Part A inspection.
\TODOfill

%======================================================================
\section{Part G: Synthesis}
%======================================================================

\subsection{When SSL provides a useful representation and when raw nonlinear classification remains competitive}
\TODOfill

\subsection{Whether the effect of SSL becomes clearer with only 10\% labels}
\TODOfill

\subsection{Stability of conclusions across the paired seeds}
\TODOfill

\subsection{Role of dataset size, dimensionality and number of classes}
\TODOfill

\subsection{Whether continuous masking improves accuracy, robustness, both or neither}
\TODOfill

\subsection{What the TF-C bearing experiment shows about transfer between operating conditions}
\TODOfill

\subsection{Implications for future multichannel industrial sensor data}
\TODOfill

%======================================================================
\section{Experimental rules: compliance notes}
%======================================================================
\begin{itemize}
    \item Official test sets preserved as held-out data: \TODOfill
    \item Test inputs not used for SSL pretraining in the required experiments: \TODOfill
    \item Normalisation, imputation and feature transforms fitted on training data only: \TODOfill
    \item 10\% label subsets stratified, indices saved: \TODOfill
    \item Same seeds, label subsets, corruptions and probe settings in paired comparisons: \TODOfill
    \item Tuning on a training-only validation split or cross-validation: \TODOfill
    \item Encoder frozen in Parts C--E; Part F follows the supplied TF-C downstream procedure: \TODOfill
    \item Configurations, software versions, runtime and hardware recorded: \TODOfill
    \item Implementation departures from the official repositories: \TODOfill
\end{itemize}

%======================================================================
\section{Limitations and conclusions}
%======================================================================
\TODOfill

%======================================================================
\appendix
\section{Optional extension: TF-C on FordA}
%======================================================================
\TODOfill

\section{Optional extension: FordB noisy-condition experiments}
\TODOfill

\section{Per-seed results}
\TODOfill

\end{document}
